\subsection{Polar decomposition}

Lemma 8. --- Let E be a separated topological vector space. Let \((\mathrm{E}_{1},\|\cdot\|_{1})\) and \((\mathrm{E}_{2},\|\cdot\|_{2})\) be dense subspaces of E endowed with Hilbert structures such that the canonical injections of \(E_{1}\) and \(E_{2}\) into E are continuous. Let F be a subspace of \(E_{1}\cap E_{2}\) , dense in \(E_{1}\) and \(E_{2}\) for these Hilbert structures. If \(\|x\|_{1}=\|x\|_{2}\) for all \(x\in F\) , then \(E_{1}=E_{2}\) and \(\|\cdot\|_{1}=\|\cdot\|_{2}\) .

Let \(x \in E_{1}\) . There exists a sequence \((x_{n})_{n \in \mathbf{N}}\) in F such that \(x_{n}\) converges to x in the Hilbert space \(E_{1}\) . Since \(\|x_{n} - x_{m}\|_{2} = \|x_{n} - x_{m}\|_{1}\) for all integers n and m, the sequence \((x_{n})\) is a Cauchy sequence in \(E_{2}\) . Let y be its limit. We have \(\|y\|_{2} = \lim\|x_{n}\|_{2} = \|x\|_{1}\) . Since the canonical injections of \(E_{1}\) and \(E_{2}\) into E are continuous, we have \(x_{n} \to x\) in E and likewise \(x_{n} \to y\) in E. Thus, it follows that x = y and \(\|x\|_{1} = \|x\|_{2}\) . In particular, \(E_{1} \subset E_{2}\) . We conclude by symmetry.

Let E be a complex Hilbert space. Let u be a positive self-adjoint partial operator on E. For every \(\alpha \in \mathbf R_{+}\) we set \(u^{\alpha} = f(u)\) , where f is the continuous map \(x \mapsto x^{\alpha}\) of \(\mathbf R_{+}\) into \(\mathbf R\). This is a positive self-adjoint partial operator (cor. 1, a) of IV, p. 273). When \(u \in \mathcal{L}(E)\) , this notation is compatible with the notation relative to the C*-algebra \(\mathcal{L}(E)\) (cf. prop. 16 of I, p. 118). If \(\alpha\) is a positive integer, the partial operator \(u^{\alpha}\) coincides with the partial operator defined by composition \(u \circ \cdots \circ u\) (cor. 1, b) of IV, p. 273).

Let \(\beta \in \mathbf{R}_{+}\) . We have \(u^{\alpha \beta} = (u^{\alpha})^{\beta}\) (cor. 4 of IV, p. 274). In particular, if \(\alpha > 0\) , then the partial operator \(u^{1 / \alpha}\) is the unique positive self-adjoint partial operator \(v\) on E such that \(v^{\alpha} = u\) .

Suppose that \(0 \leqslant \alpha \leqslant \beta\) . We then have \(\operatorname{dom}(u^{\beta}) \subset \operatorname{dom}(u^{\alpha})\) : indeed, for every real number \(x \geqslant 0\) , we have \(x^{\alpha} \leqslant 1 + x^{\beta}\) , and the result then follows from the definition of the domain of \(u^{\alpha}\) (cf. prop. 3 of IV, p. 271).

Let u be a closed densely defined operator from E into a complex Hilbert space F. The partial operator \(u^{*}u\) is self-adjoint and positive (prop. 12 of IV, p. 241). We denote \(|u| = (u^{*}u)^{1/2}\) .

PROPOSITION 12. --- Let u be a closed densely defined operator from E into a complex Hilbert space F.

\begin{enumerate}
\def\labelenumi{\alph{enumi})}
\item
  The domain of the positive self-adjoint partial operator \(|u|\) coincides with the domain of \(u\) ;
\item
  There exists a unique partially isometric linear map \(j\) from E into F such that \(u = j|u|\) and \(\operatorname{Ker}(j) = \operatorname{Ker}(u)\) ;
\item
  Let \(u_{1}\) be a positive self-adjoint operator on E and \(j_{1}\) a partially isometric linear map from E into F such that \(u = j_{1}u_{1}\) and \(\operatorname{Ker}(j_{1}) = \operatorname{Ker}(u_{1})\) . We then have \(u_{1} = |u|\) and \(j_{1} = j\) .
\end{enumerate}

The domain of \(u^{*}u\) is contained in \(\operatorname{dom}(u)\) and in \(\operatorname{dom}(|u|)\) . It is dense in the Hilbert spaces \(E_{u}\) (loc. cit.) and \(E_{|u|}\) (cor. 1, e) of IV, p. 273) and, for every \(x \in \operatorname{dom}(u^{*}u)\) , we have

\[
\begin{array}{c} \| x \| _ {u} ^ {2} = \| x \| ^ {2} + \langle u (x) | u (x) \rangle = \| x \| ^ {2} + \langle x | (u ^ {*} u) (x) \rangle \\ = \| x \| ^ {2} + \langle | u | (x) | | u | (x) \rangle = \| x \| _ {| u |} ^ {2}. \end{array}\tag{16}
\]

Consequently, the Hilbert spaces \(E_{u}\) and \(E_{|u|}\) are equal (lemma 8), hence \(\text{dom}(u) = \text{dom}(|u|)\) and \(\|u(x)\| = \||u|(x)\|\) for all \(x \in \text{dom}(u)\) .

Formula (16) implies that \(\operatorname{Ker}(u) = \operatorname{Ker}(|u|)\) and that there exists a unique isometric linear map \(v\) from \(\overline{\operatorname{Im}(|u|)}\) onto \(\overline{\operatorname{Im}(u)}\) satisfying \(v(|u|(x)) = u(x)\) for all \(x \in \operatorname{dom}(|u|)\) . Since \(|u|\) is self-adjoint, we have \(\operatorname{Im}(|u|)^{\circ} = \operatorname{Ker}(|u|)\) (prop. 7, c) of IV, p. 236). There therefore exists a unique partial isometry \(j\) from E into F which extends \(v\) and vanishes on \(\operatorname{Ker}(|u|) = \operatorname{Ker}(u)\) . Since \(\operatorname{E} = \operatorname{Ker}(u) \oplus \overline{\operatorname{Im}(|u|)}\) , this map is the unique partially isometric map such that \(u = j|u|\) and \(\operatorname{Ker}(j) = \operatorname{Ker}(u)\) .

Let us prove c). We have \(u_{1}j_{1}^{*}j_{1}u_{1}\subset(j_{1}u_{1})^{*}j_{1}u_{1}=u^{*}u\) . The linear map \(j_{1}^{*}j_{1}\) is the orthogonal projector with kernel \(\mathrm{Ker}(j_{1})=\mathrm{Ker}(u_{1})\) (EVT, V, p. 41, prop. 5 (ii)) and therefore with image \(\mathrm{Ker}(u_{1})^{\circ}=\overline{\mathrm{Im}(u_{1})}\) (prop. 7, c) of IV, p. 236). Consequently, \(u_{1}^{2}\subset u^{*}u\) , whence \(u_{1}^{2}=u^{*}u\) since these two operators are self-adjoint. Thus, it follows that \(u_{1}=(u^{*}u)^{1/2}\) and the uniqueness assertion of b) ultimately shows that \(j_{1}=j\) .

DEFINITION 7. --- Let u be a closed operator with dense domain from E into a complex Hilbert space F. The pair \((j, |u|)\) determined by prop. 12 is called the polar decomposition of u.

Suppose that \(u \in \mathcal{L}(\mathrm{E};\mathrm{F})\) . Its polar decomposition in the sense of this definition coincides with that of Definition 4 in I, p. 140.

